
% ================================================================
% Linear Algebra Textbook Preamble and Macros
% ================================================================
%
% Suggested main.tex:
%
% \documentclass[11pt,oneside]{book}
% 
% ================================================================
% Linear Algebra Textbook Preamble and Macros
% ================================================================
%
% Suggested main.tex:
%
% \documentclass[11pt,oneside]{book}
% 
% ================================================================
% Linear Algebra Textbook Preamble and Macros
% ================================================================
%
% Suggested main.tex:
%
% \documentclass[11pt,oneside]{book}
% 
% ================================================================
% Linear Algebra Textbook Preamble and Macros
% ================================================================
%
% Suggested main.tex:
%
% \documentclass[11pt,oneside]{book}
% \input{preamble}
%
% \begin{document}
% \frontmatter
% \maketitle
% \tableofcontents
%
% \mainmatter
% \chapter{Introduction}
% ...
%
% \end{document}
%
% ================================================================


% ================================================================
% BASIC PACKAGES
% ================================================================

\usepackage[utf8]{inputenc}
\usepackage[T1]{fontenc}
\usepackage{lmodern}
\usepackage{microtype}

\usepackage[margin=1in]{geometry}

\usepackage{amsmath, amssymb, amsthm, mathtools}
\usepackage{graphicx}
\usepackage{xcolor}
\usepackage{booktabs}
\usepackage{array}
\usepackage{enumitem}
\usepackage{comment}
\usepackage{xparse}
\usepackage{etoolbox}


% ================================================================
% FIGURES AND DRAWINGS
% ================================================================

\usepackage{tikz}

\usetikzlibrary{
  arrows.meta,
  calc,
  intersections,
  positioning,
  decorations.pathreplacing,
  angles,
  quotes
}



\usepackage{pgfplots}
\pgfplotsset{compat=1.18}

\definecolor{notespink}{RGB}{210,45,105}
\definecolor{notespurple}{RGB}{95,45,170}


% ================================================================
% BOXES
% ================================================================

\usepackage[most]{tcolorbox}


% ================================================================
% HYPERLINKS AND REFERENCES
% ================================================================

\usepackage{hyperref}

\hypersetup{
  colorlinks=true,
  linkcolor=blue!60!black,
  citecolor=blue!60!black,
  urlcolor=blue!70!black,
  pdftitle={Linear Algebra},
  pdfauthor={},
  pdfsubject={Linear Algebra},
  pdfkeywords={linear algebra, vectors, matrices}
}

\usepackage[nameinlink]{cleveref}


% ================================================================
% PAGE STYLE
% ================================================================

\usepackage{fancyhdr}

\pagestyle{fancy}
\fancyhf{}
\fancyhead[L]{Linear Algebra}
\fancyhead[R]{\leftmark}
\fancyfoot[C]{\thepage}

\setlength{\headheight}{14pt}


% ================================================================
% PARAGRAPH AND LIST FORMATTING
% ================================================================

\setlength{\parindent}{0pt}
\setlength{\parskip}{0.65em}

\setlist[itemize]{
  leftmargin=2em,
  itemsep=0.25em,
  topsep=0.25em
}

\setlist[enumerate]{
  leftmargin=2em,
  itemsep=0.35em,
  topsep=0.35em
}


% ================================================================
% NUMBERING
% ================================================================

\numberwithin{equation}{chapter}


% ================================================================
% COMMON NUMBER SYSTEMS
% ================================================================

\newcommand{\R}{\mathbb R}
\newcommand{\Rone}{\mathbb R^1}
\newcommand{\Rtwo}{\mathbb R^2}
\newcommand{\Rthree}{\mathbb R^3}
\newcommand{\Rn}{\mathbb R^n}
\newcommand{\Rm}{\mathbb R^m}

\newcommand{\N}{\mathbb N}
\newcommand{\Z}{\mathbb Z}
\newcommand{\Q}{\mathbb Q}
\newcommand{\C}{\mathbb C}


% ================================================================
% SET NOTATION
% ================================================================

% Example:
%   \Set{x \in \R}{x > 0}
%
% Produces:
%   { x in R : x > 0 }

\newcommand{\Set}[2]{\left\{\, #1 \;:\; #2 \,\right\}}

% Alternative with vertical bar:
%   \Setbar{x \in \R}{x > 0}

\newcommand{\Setbar}[2]{\left\{\, #1 \;\middle|\; #2 \,\right\}}

\newcommand{\st}{\;:\;}
\newcommand{\suchthat}{\;\middle|\;}


% ================================================================
% VECTOR NOTATION
% ================================================================
%
% We use arrow notation throughout:
%
%   \vec v, \vec u, \vec x
%
% We do not use boldface for vectors.
% The standard LaTeX \vec command is left unchanged.

\newcommand{\vzero}{\vec 0}

\newcommand{\va}{\vec a}
\newcommand{\vb}{\vec b}
\newcommand{\vc}{\vec c}
\newcommand{\vu}{\vec u}
\newcommand{\vv}{\vec v}
\newcommand{\vw}{\vec w}
\newcommand{\vx}{\vec x}
\newcommand{\vy}{\vec y}
\newcommand{\vz}{\vec z}
\newcommand{\vp}{\vec p}
\newcommand{\vq}{\vec q}

% Standard basis vectors.

\newcommand{\eone}{\vec e_1}
\newcommand{\etwo}{\vec e_2}
\newcommand{\ethree}{\vec e_3}
\newcommand{\ei}{\vec e_i}
\newcommand{\ej}{\vec e_j}

% Optional physics/engineering notation.
% Main notes should primarily use \vec e_1, \vec e_2, \vec e_3.

\newcommand{\ihat}{\hat{\imath}}
\newcommand{\jhat}{\hat{\jmath}}
\newcommand{\khat}{\hat{k}}


% ================================================================
% COLUMN VECTORS AND MATRICES
% ================================================================
%
% We emphasize column vectors.
%
% Examples:
%
%   \twovec{2}{3}
%   \threevec{1}{0}{-4}
%   \colvec{x_1 \\ x_2 \\ \vdots \\ x_n}
%   \mat{1 & 2 \\ 3 & 4}

\newcommand{\colvec}[1]{\begin{bmatrix} #1 \end{bmatrix}}

\newcommand{\twovec}[2]{
\begin{bmatrix}
#1\\
#2
\end{bmatrix}
}

\newcommand{\threevec}[3]{
\begin{bmatrix}
#1\\
#2\\
#3
\end{bmatrix}
}

\newcommand{\fourvec}[4]{
\begin{bmatrix}
#1\\
#2\\
#3\\
#4
\end{bmatrix}
}

\newcommand{\mat}[1]{\begin{bmatrix} #1 \end{bmatrix}}

\newcommand{\smallmat}[1]{%
\left[\begin{smallmatrix} #1 \end{smallmatrix}\right]%
}

% Augmented matrix.
%
% Example:
%
% \augmat{cc|c}{
% 1 & 2 & 3\\
% 4 & 5 & 6
% }

\newcommand{\augmat}[2]{
\left[
\begin{array}{#1}
#2
\end{array}
\right]
}


% ================================================================
% INNER PRODUCTS, NORMS, DISTANCE
% ================================================================

\newcommand{\dotp}[2]{#1 \cdot #2}
\newcommand{\inner}[2]{\left\langle #1, #2 \right\rangle}

\newcommand{\norm}[1]{\left\lVert #1 \right\rVert}
\newcommand{\abs}[1]{\left| #1 \right|}
\newcommand{\dist}[2]{\operatorname{dist}\!\left(#1,#2\right)}


% ================================================================
% LINEAR ALGEBRA OPERATORS
% ================================================================

\DeclareMathOperator{\Span}{span}
\DeclareMathOperator{\Col}{Col}
\DeclareMathOperator{\Nul}{Nul}
\DeclareMathOperator{\Row}{Row}
\DeclareMathOperator{\rank}{rank}
\DeclareMathOperator{\nullity}{nullity}
\DeclareMathOperator{\im}{im}
\DeclareMathOperator{\range}{range}
\DeclareMathOperator{\projop}{proj}
\DeclareMathOperator{\diag}{diag}
\DeclareMathOperator{\tr}{tr}

% Projection notation.
%
% Example:
%   \proj{\vec u}{\vec v}
%
% Means:
%   projection of \vec v onto the direction/subspace \vec u.

\newcommand{\proj}[2]{\projop_{#1}\!\left(#2\right)}

% Matrix transpose and inverse.
%
% Examples:
%   A\T
%   A\inv

\newcommand{\T}{^{\mathsf T}}
\newcommand{\inv}{^{-1}}

% Coordinates relative to a basis.
%
% Example:
%   \coords{\vec v}{B}
%
% Produces:
%   [v]_B

\newcommand{\coords}[2]{\left[#1\right]_{#2}}


% ================================================================
% TEXT EMPHASIS MACROS
% ================================================================

\newcommand{\defterm}[1]{\textbf{#1}}

\newcommand{\important}[1]{%
\begin{center}
\fbox{\parbox{0.85\textwidth}{#1}}
\end{center}
}

\newcommand{\bigidea}[1]{%
\begin{center}
\begin{tcolorbox}[
  width=0.88\textwidth,
  colback=yellow!8,
  colframe=yellow!50!black,
  boxrule=0.8pt,
  arc=2mm
]
#1
\end{tcolorbox}
\end{center}
}


% ================================================================
% THEOREM-LIKE ENVIRONMENTS
% ================================================================

\theoremstyle{definition}

\newtheorem{definition}{Definition}[chapter]
\newtheorem{example}{Example}[chapter]
\newtheorem{exerciseinner}{Exercise}[chapter]

\theoremstyle{plain}

\newtheorem{theorem}{Theorem}[chapter]
\newtheorem{proposition}{Proposition}[chapter]
\newtheorem{fact}{Fact}[chapter]

\theoremstyle{remark}

\newtheorem{remark}{Remark}[chapter]


% ================================================================
% BOX STYLING FOR THEOREM-LIKE ENVIRONMENTS
% ================================================================

\tcolorboxenvironment{definition}{
  enhanced,
  breakable,
  colback=blue!3,
  colframe=blue!55!black,
  boxrule=0.7pt,
  arc=2mm,
  left=1mm,
  right=1mm,
  top=1mm,
  bottom=1mm
}

\tcolorboxenvironment{example}{
  enhanced,
  breakable,
  colback=green!3,
  colframe=green!45!black,
  boxrule=0.7pt,
  arc=2mm,
  left=1mm,
  right=1mm,
  top=1mm,
  bottom=1mm
}

\tcolorboxenvironment{theorem}{
  enhanced,
  breakable,
  colback=purple!3,
  colframe=purple!55!black,
  boxrule=0.7pt,
  arc=2mm,
  left=1mm,
  right=1mm,
  top=1mm,
  bottom=1mm
}

\tcolorboxenvironment{proposition}{
  enhanced,
  breakable,
  colback=purple!3,
  colframe=purple!45!black,
  boxrule=0.7pt,
  arc=2mm,
  left=1mm,
  right=1mm,
  top=1mm,
  bottom=1mm
}

\tcolorboxenvironment{fact}{
  enhanced,
  breakable,
  colback=orange!5,
  colframe=orange!60!black,
  boxrule=0.7pt,
  arc=2mm,
  left=1mm,
  right=1mm,
  top=1mm,
  bottom=1mm
}

\tcolorboxenvironment{remark}{
  enhanced,
  breakable,
  colback=gray!5,
  colframe=gray!50!black,
  boxrule=0.5pt,
  arc=2mm,
  left=1mm,
  right=1mm,
  top=1mm,
  bottom=1mm
}


% ================================================================
% EXERCISES AND SOLUTIONS
% ================================================================
%
% Exercise:
%
% \begin{exercise}
% ...
% \end{exercise}
%
% Exercise with title:
%
% \begin{exercise}[Interpreting a line]
% ...
% \end{exercise}

\NewDocumentEnvironment{exercise}{O{}}
{
  \ifstrempty{#1}
    {\begin{exerciseinner}}
    {\begin{exerciseinner}[#1]}
}
{
  \end{exerciseinner}
}

% Solutions are hidden by default.
% To show solutions, change \showsolutionstrue below.

\newif\ifshowsolutions
\showsolutionsfalse
% \showsolutionstrue

\ifshowsolutions
  \NewDocumentEnvironment{solution}{}
  {
    \begin{proof}[Solution]
  }
  {
    \end{proof}
  }
\else
  \excludecomment{solution}
\fi


% ================================================================
% CHAPTER-LEVEL STRUCTURE BOXES
% ================================================================

\newenvironment{chapteroverview}
{
  \begin{tcolorbox}[
    enhanced,
    breakable,
    colback=cyan!4,
    colframe=cyan!50!black,
    title=\textbf{Chapter Overview},
    boxrule=0.7pt,
    arc=2mm
  ]
}
{
  \end{tcolorbox}
}

\newenvironment{learningobjectives}
{
  \begin{tcolorbox}[
    enhanced,
    breakable,
    colback=cyan!4,
    colframe=cyan!50!black,
    title=\textbf{Learning Objectives},
    boxrule=0.7pt,
    arc=2mm
  ]
  \begin{itemize}
}
{
  \end{itemize}
  \end{tcolorbox}
}

\newenvironment{summarybox}
{
  \begin{tcolorbox}[
    enhanced,
    breakable,
    colback=yellow!5,
    colframe=yellow!50!black,
    title=\textbf{Summary},
    boxrule=0.7pt,
    arc=2mm
  ]
  \begin{itemize}
}
{
  \end{itemize}
  \end{tcolorbox}
}

\newenvironment{checkyourunderstanding}
{
  \begin{tcolorbox}[
    enhanced,
    breakable,
    colback=red!3,
    colframe=red!50!black,
    title=\textbf{Check Your Understanding},
    boxrule=0.7pt,
    arc=2mm
  ]
  \begin{enumerate}
}
{
  \end{enumerate}
  \end{tcolorbox}
}


% ================================================================
% OPTIONAL SECTION-LEVEL BOXES
% ================================================================

\newenvironment{motivation}
{
  \begin{tcolorbox}[
    enhanced,
    breakable,
    colback=green!4,
    colframe=green!45!black,
    title=\textbf{Motivation},
    boxrule=0.7pt,
    arc=2mm
  ]
}
{
  \end{tcolorbox}
}

\newenvironment{warningbox}
{
  \begin{tcolorbox}[
    enhanced,
    breakable,
    colback=red!4,
    colframe=red!60!black,
    title=\textbf{Warning},
    boxrule=0.7pt,
    arc=2mm
  ]
}
{
  \end{tcolorbox}
}

\newenvironment{computationalnote}
{
  \begin{tcolorbox}[
    enhanced,
    breakable,
    colback=gray!6,
    colframe=gray!60!black,
    title=\textbf{Computational Note},
    boxrule=0.7pt,
    arc=2mm
  ]
}
{
  \end{tcolorbox}
}


% ================================================================
% TIKZ STYLES
% ================================================================

\tikzset{
  vector/.style={
    -{Latex[length=3mm]},
    thick,
    blue!70!black
  },
  axis/.style={
    -{Latex[length=2.5mm]},
    gray!70
  },
  point/.style={
    circle,
    fill=black,
    inner sep=1.5pt
  },
  dashedline/.style={
    dashed,
    gray!70
  }
}

% Simple coordinate axes for \R^2 sketches.
%
% Usage inside tikzpicture:
%
% \drawRtwoaxes{-3}{3}{-2}{2}
%
% Arguments:
%   xmin, xmax, ymin, ymax

\NewDocumentCommand{\drawRtwoaxes}{m m m m}
{
  \draw[axis] (#1,0) -- (#2,0) node[right] {$x$};
  \draw[axis] (0,#3) -- (0,#4) node[above] {$y$};
}


% ================================================================
% CLEVEREF NAMES
% ================================================================

\crefname{definition}{Definition}{Definitions}
\Crefname{definition}{Definition}{Definitions}

\crefname{example}{Example}{Examples}
\Crefname{example}{Example}{Examples}

\crefname{theorem}{Theorem}{Theorems}
\Crefname{theorem}{Theorem}{Theorems}

\crefname{proposition}{Proposition}{Propositions}
\Crefname{proposition}{Proposition}{Propositions}

\crefname{fact}{Fact}{Facts}
\Crefname{fact}{Fact}{Facts}

\crefname{remark}{Remark}{Remarks}
\Crefname{remark}{Remark}{Remarks}

\crefname{exerciseinner}{Exercise}{Exercises}
\Crefname{exerciseinner}{Exercise}{Exercises}


% ================================================================
% END OF PREAMBLE
% ================================================================

%
% \begin{document}
% \frontmatter
% \maketitle
% \tableofcontents
%
% \mainmatter
% \chapter{Introduction}
% ...
%
% \end{document}
%
% ================================================================


% ================================================================
% BASIC PACKAGES
% ================================================================

\usepackage[utf8]{inputenc}
\usepackage[T1]{fontenc}
\usepackage{lmodern}
\usepackage{microtype}

\usepackage[margin=1in]{geometry}

\usepackage{amsmath, amssymb, amsthm, mathtools}
\usepackage{graphicx}
\usepackage{xcolor}
\usepackage{booktabs}
\usepackage{array}
\usepackage{enumitem}
\usepackage{comment}
\usepackage{xparse}
\usepackage{etoolbox}


% ================================================================
% FIGURES AND DRAWINGS
% ================================================================

\usepackage{tikz}

\usetikzlibrary{
  arrows.meta,
  calc,
  intersections,
  positioning,
  decorations.pathreplacing,
  angles,
  quotes
}



\usepackage{pgfplots}
\pgfplotsset{compat=1.18}

\definecolor{notespink}{RGB}{210,45,105}
\definecolor{notespurple}{RGB}{95,45,170}


% ================================================================
% BOXES
% ================================================================

\usepackage[most]{tcolorbox}


% ================================================================
% HYPERLINKS AND REFERENCES
% ================================================================

\usepackage{hyperref}

\hypersetup{
  colorlinks=true,
  linkcolor=blue!60!black,
  citecolor=blue!60!black,
  urlcolor=blue!70!black,
  pdftitle={Linear Algebra},
  pdfauthor={},
  pdfsubject={Linear Algebra},
  pdfkeywords={linear algebra, vectors, matrices}
}

\usepackage[nameinlink]{cleveref}


% ================================================================
% PAGE STYLE
% ================================================================

\usepackage{fancyhdr}

\pagestyle{fancy}
\fancyhf{}
\fancyhead[L]{Linear Algebra}
\fancyhead[R]{\leftmark}
\fancyfoot[C]{\thepage}

\setlength{\headheight}{14pt}


% ================================================================
% PARAGRAPH AND LIST FORMATTING
% ================================================================

\setlength{\parindent}{0pt}
\setlength{\parskip}{0.65em}

\setlist[itemize]{
  leftmargin=2em,
  itemsep=0.25em,
  topsep=0.25em
}

\setlist[enumerate]{
  leftmargin=2em,
  itemsep=0.35em,
  topsep=0.35em
}


% ================================================================
% NUMBERING
% ================================================================

\numberwithin{equation}{chapter}


% ================================================================
% COMMON NUMBER SYSTEMS
% ================================================================

\newcommand{\R}{\mathbb R}
\newcommand{\Rone}{\mathbb R^1}
\newcommand{\Rtwo}{\mathbb R^2}
\newcommand{\Rthree}{\mathbb R^3}
\newcommand{\Rn}{\mathbb R^n}
\newcommand{\Rm}{\mathbb R^m}

\newcommand{\N}{\mathbb N}
\newcommand{\Z}{\mathbb Z}
\newcommand{\Q}{\mathbb Q}
\newcommand{\C}{\mathbb C}


% ================================================================
% SET NOTATION
% ================================================================

% Example:
%   \Set{x \in \R}{x > 0}
%
% Produces:
%   { x in R : x > 0 }

\newcommand{\Set}[2]{\left\{\, #1 \;:\; #2 \,\right\}}

% Alternative with vertical bar:
%   \Setbar{x \in \R}{x > 0}

\newcommand{\Setbar}[2]{\left\{\, #1 \;\middle|\; #2 \,\right\}}

\newcommand{\st}{\;:\;}
\newcommand{\suchthat}{\;\middle|\;}


% ================================================================
% VECTOR NOTATION
% ================================================================
%
% We use arrow notation throughout:
%
%   \vec v, \vec u, \vec x
%
% We do not use boldface for vectors.
% The standard LaTeX \vec command is left unchanged.

\newcommand{\vzero}{\vec 0}

\newcommand{\va}{\vec a}
\newcommand{\vb}{\vec b}
\newcommand{\vc}{\vec c}
\newcommand{\vu}{\vec u}
\newcommand{\vv}{\vec v}
\newcommand{\vw}{\vec w}
\newcommand{\vx}{\vec x}
\newcommand{\vy}{\vec y}
\newcommand{\vz}{\vec z}
\newcommand{\vp}{\vec p}
\newcommand{\vq}{\vec q}

% Standard basis vectors.

\newcommand{\eone}{\vec e_1}
\newcommand{\etwo}{\vec e_2}
\newcommand{\ethree}{\vec e_3}
\newcommand{\ei}{\vec e_i}
\newcommand{\ej}{\vec e_j}

% Optional physics/engineering notation.
% Main notes should primarily use \vec e_1, \vec e_2, \vec e_3.

\newcommand{\ihat}{\hat{\imath}}
\newcommand{\jhat}{\hat{\jmath}}
\newcommand{\khat}{\hat{k}}


% ================================================================
% COLUMN VECTORS AND MATRICES
% ================================================================
%
% We emphasize column vectors.
%
% Examples:
%
%   \twovec{2}{3}
%   \threevec{1}{0}{-4}
%   \colvec{x_1 \\ x_2 \\ \vdots \\ x_n}
%   \mat{1 & 2 \\ 3 & 4}

\newcommand{\colvec}[1]{\begin{bmatrix} #1 \end{bmatrix}}

\newcommand{\twovec}[2]{
\begin{bmatrix}
#1\\
#2
\end{bmatrix}
}

\newcommand{\threevec}[3]{
\begin{bmatrix}
#1\\
#2\\
#3
\end{bmatrix}
}

\newcommand{\fourvec}[4]{
\begin{bmatrix}
#1\\
#2\\
#3\\
#4
\end{bmatrix}
}

\newcommand{\mat}[1]{\begin{bmatrix} #1 \end{bmatrix}}

\newcommand{\smallmat}[1]{%
\left[\begin{smallmatrix} #1 \end{smallmatrix}\right]%
}

% Augmented matrix.
%
% Example:
%
% \augmat{cc|c}{
% 1 & 2 & 3\\
% 4 & 5 & 6
% }

\newcommand{\augmat}[2]{
\left[
\begin{array}{#1}
#2
\end{array}
\right]
}


% ================================================================
% INNER PRODUCTS, NORMS, DISTANCE
% ================================================================

\newcommand{\dotp}[2]{#1 \cdot #2}
\newcommand{\inner}[2]{\left\langle #1, #2 \right\rangle}

\newcommand{\norm}[1]{\left\lVert #1 \right\rVert}
\newcommand{\abs}[1]{\left| #1 \right|}
\newcommand{\dist}[2]{\operatorname{dist}\!\left(#1,#2\right)}


% ================================================================
% LINEAR ALGEBRA OPERATORS
% ================================================================

\DeclareMathOperator{\Span}{span}
\DeclareMathOperator{\Col}{Col}
\DeclareMathOperator{\Nul}{Nul}
\DeclareMathOperator{\Row}{Row}
\DeclareMathOperator{\rank}{rank}
\DeclareMathOperator{\nullity}{nullity}
\DeclareMathOperator{\im}{im}
\DeclareMathOperator{\range}{range}
\DeclareMathOperator{\projop}{proj}
\DeclareMathOperator{\diag}{diag}
\DeclareMathOperator{\tr}{tr}

% Projection notation.
%
% Example:
%   \proj{\vec u}{\vec v}
%
% Means:
%   projection of \vec v onto the direction/subspace \vec u.

\newcommand{\proj}[2]{\projop_{#1}\!\left(#2\right)}

% Matrix transpose and inverse.
%
% Examples:
%   A\T
%   A\inv

\newcommand{\T}{^{\mathsf T}}
\newcommand{\inv}{^{-1}}

% Coordinates relative to a basis.
%
% Example:
%   \coords{\vec v}{B}
%
% Produces:
%   [v]_B

\newcommand{\coords}[2]{\left[#1\right]_{#2}}


% ================================================================
% TEXT EMPHASIS MACROS
% ================================================================

\newcommand{\defterm}[1]{\textbf{#1}}

\newcommand{\important}[1]{%
\begin{center}
\fbox{\parbox{0.85\textwidth}{#1}}
\end{center}
}

\newcommand{\bigidea}[1]{%
\begin{center}
\begin{tcolorbox}[
  width=0.88\textwidth,
  colback=yellow!8,
  colframe=yellow!50!black,
  boxrule=0.8pt,
  arc=2mm
]
#1
\end{tcolorbox}
\end{center}
}


% ================================================================
% THEOREM-LIKE ENVIRONMENTS
% ================================================================

\theoremstyle{definition}

\newtheorem{definition}{Definition}[chapter]
\newtheorem{example}{Example}[chapter]
\newtheorem{exerciseinner}{Exercise}[chapter]

\theoremstyle{plain}

\newtheorem{theorem}{Theorem}[chapter]
\newtheorem{proposition}{Proposition}[chapter]
\newtheorem{fact}{Fact}[chapter]

\theoremstyle{remark}

\newtheorem{remark}{Remark}[chapter]


% ================================================================
% BOX STYLING FOR THEOREM-LIKE ENVIRONMENTS
% ================================================================

\tcolorboxenvironment{definition}{
  enhanced,
  breakable,
  colback=blue!3,
  colframe=blue!55!black,
  boxrule=0.7pt,
  arc=2mm,
  left=1mm,
  right=1mm,
  top=1mm,
  bottom=1mm
}

\tcolorboxenvironment{example}{
  enhanced,
  breakable,
  colback=green!3,
  colframe=green!45!black,
  boxrule=0.7pt,
  arc=2mm,
  left=1mm,
  right=1mm,
  top=1mm,
  bottom=1mm
}

\tcolorboxenvironment{theorem}{
  enhanced,
  breakable,
  colback=purple!3,
  colframe=purple!55!black,
  boxrule=0.7pt,
  arc=2mm,
  left=1mm,
  right=1mm,
  top=1mm,
  bottom=1mm
}

\tcolorboxenvironment{proposition}{
  enhanced,
  breakable,
  colback=purple!3,
  colframe=purple!45!black,
  boxrule=0.7pt,
  arc=2mm,
  left=1mm,
  right=1mm,
  top=1mm,
  bottom=1mm
}

\tcolorboxenvironment{fact}{
  enhanced,
  breakable,
  colback=orange!5,
  colframe=orange!60!black,
  boxrule=0.7pt,
  arc=2mm,
  left=1mm,
  right=1mm,
  top=1mm,
  bottom=1mm
}

\tcolorboxenvironment{remark}{
  enhanced,
  breakable,
  colback=gray!5,
  colframe=gray!50!black,
  boxrule=0.5pt,
  arc=2mm,
  left=1mm,
  right=1mm,
  top=1mm,
  bottom=1mm
}


% ================================================================
% EXERCISES AND SOLUTIONS
% ================================================================
%
% Exercise:
%
% \begin{exercise}
% ...
% \end{exercise}
%
% Exercise with title:
%
% \begin{exercise}[Interpreting a line]
% ...
% \end{exercise}

\NewDocumentEnvironment{exercise}{O{}}
{
  \ifstrempty{#1}
    {\begin{exerciseinner}}
    {\begin{exerciseinner}[#1]}
}
{
  \end{exerciseinner}
}

% Solutions are hidden by default.
% To show solutions, change \showsolutionstrue below.

\newif\ifshowsolutions
\showsolutionsfalse
% \showsolutionstrue

\ifshowsolutions
  \NewDocumentEnvironment{solution}{}
  {
    \begin{proof}[Solution]
  }
  {
    \end{proof}
  }
\else
  \excludecomment{solution}
\fi


% ================================================================
% CHAPTER-LEVEL STRUCTURE BOXES
% ================================================================

\newenvironment{chapteroverview}
{
  \begin{tcolorbox}[
    enhanced,
    breakable,
    colback=cyan!4,
    colframe=cyan!50!black,
    title=\textbf{Chapter Overview},
    boxrule=0.7pt,
    arc=2mm
  ]
}
{
  \end{tcolorbox}
}

\newenvironment{learningobjectives}
{
  \begin{tcolorbox}[
    enhanced,
    breakable,
    colback=cyan!4,
    colframe=cyan!50!black,
    title=\textbf{Learning Objectives},
    boxrule=0.7pt,
    arc=2mm
  ]
  \begin{itemize}
}
{
  \end{itemize}
  \end{tcolorbox}
}

\newenvironment{summarybox}
{
  \begin{tcolorbox}[
    enhanced,
    breakable,
    colback=yellow!5,
    colframe=yellow!50!black,
    title=\textbf{Summary},
    boxrule=0.7pt,
    arc=2mm
  ]
  \begin{itemize}
}
{
  \end{itemize}
  \end{tcolorbox}
}

\newenvironment{checkyourunderstanding}
{
  \begin{tcolorbox}[
    enhanced,
    breakable,
    colback=red!3,
    colframe=red!50!black,
    title=\textbf{Check Your Understanding},
    boxrule=0.7pt,
    arc=2mm
  ]
  \begin{enumerate}
}
{
  \end{enumerate}
  \end{tcolorbox}
}


% ================================================================
% OPTIONAL SECTION-LEVEL BOXES
% ================================================================

\newenvironment{motivation}
{
  \begin{tcolorbox}[
    enhanced,
    breakable,
    colback=green!4,
    colframe=green!45!black,
    title=\textbf{Motivation},
    boxrule=0.7pt,
    arc=2mm
  ]
}
{
  \end{tcolorbox}
}

\newenvironment{warningbox}
{
  \begin{tcolorbox}[
    enhanced,
    breakable,
    colback=red!4,
    colframe=red!60!black,
    title=\textbf{Warning},
    boxrule=0.7pt,
    arc=2mm
  ]
}
{
  \end{tcolorbox}
}

\newenvironment{computationalnote}
{
  \begin{tcolorbox}[
    enhanced,
    breakable,
    colback=gray!6,
    colframe=gray!60!black,
    title=\textbf{Computational Note},
    boxrule=0.7pt,
    arc=2mm
  ]
}
{
  \end{tcolorbox}
}


% ================================================================
% TIKZ STYLES
% ================================================================

\tikzset{
  vector/.style={
    -{Latex[length=3mm]},
    thick,
    blue!70!black
  },
  axis/.style={
    -{Latex[length=2.5mm]},
    gray!70
  },
  point/.style={
    circle,
    fill=black,
    inner sep=1.5pt
  },
  dashedline/.style={
    dashed,
    gray!70
  }
}

% Simple coordinate axes for \R^2 sketches.
%
% Usage inside tikzpicture:
%
% \drawRtwoaxes{-3}{3}{-2}{2}
%
% Arguments:
%   xmin, xmax, ymin, ymax

\NewDocumentCommand{\drawRtwoaxes}{m m m m}
{
  \draw[axis] (#1,0) -- (#2,0) node[right] {$x$};
  \draw[axis] (0,#3) -- (0,#4) node[above] {$y$};
}


% ================================================================
% CLEVEREF NAMES
% ================================================================

\crefname{definition}{Definition}{Definitions}
\Crefname{definition}{Definition}{Definitions}

\crefname{example}{Example}{Examples}
\Crefname{example}{Example}{Examples}

\crefname{theorem}{Theorem}{Theorems}
\Crefname{theorem}{Theorem}{Theorems}

\crefname{proposition}{Proposition}{Propositions}
\Crefname{proposition}{Proposition}{Propositions}

\crefname{fact}{Fact}{Facts}
\Crefname{fact}{Fact}{Facts}

\crefname{remark}{Remark}{Remarks}
\Crefname{remark}{Remark}{Remarks}

\crefname{exerciseinner}{Exercise}{Exercises}
\Crefname{exerciseinner}{Exercise}{Exercises}


% ================================================================
% END OF PREAMBLE
% ================================================================

%
% \begin{document}
% \frontmatter
% \maketitle
% \tableofcontents
%
% \mainmatter
% \chapter{Introduction}
% ...
%
% \end{document}
%
% ================================================================


% ================================================================
% BASIC PACKAGES
% ================================================================

\usepackage[utf8]{inputenc}
\usepackage[T1]{fontenc}
\usepackage{lmodern}
\usepackage{microtype}

\usepackage[margin=1in]{geometry}

\usepackage{amsmath, amssymb, amsthm, mathtools}
\usepackage{graphicx}
\usepackage{xcolor}
\usepackage{booktabs}
\usepackage{array}
\usepackage{enumitem}
\usepackage{comment}
\usepackage{xparse}
\usepackage{etoolbox}


% ================================================================
% FIGURES AND DRAWINGS
% ================================================================

\usepackage{tikz}

\usetikzlibrary{
  arrows.meta,
  calc,
  intersections,
  positioning,
  decorations.pathreplacing,
  angles,
  quotes
}



\usepackage{pgfplots}
\pgfplotsset{compat=1.18}

\definecolor{notespink}{RGB}{210,45,105}
\definecolor{notespurple}{RGB}{95,45,170}


% ================================================================
% BOXES
% ================================================================

\usepackage[most]{tcolorbox}


% ================================================================
% HYPERLINKS AND REFERENCES
% ================================================================

\usepackage{hyperref}

\hypersetup{
  colorlinks=true,
  linkcolor=blue!60!black,
  citecolor=blue!60!black,
  urlcolor=blue!70!black,
  pdftitle={Linear Algebra},
  pdfauthor={},
  pdfsubject={Linear Algebra},
  pdfkeywords={linear algebra, vectors, matrices}
}

\usepackage[nameinlink]{cleveref}


% ================================================================
% PAGE STYLE
% ================================================================

\usepackage{fancyhdr}

\pagestyle{fancy}
\fancyhf{}
\fancyhead[L]{Linear Algebra}
\fancyhead[R]{\leftmark}
\fancyfoot[C]{\thepage}

\setlength{\headheight}{14pt}


% ================================================================
% PARAGRAPH AND LIST FORMATTING
% ================================================================

\setlength{\parindent}{0pt}
\setlength{\parskip}{0.65em}

\setlist[itemize]{
  leftmargin=2em,
  itemsep=0.25em,
  topsep=0.25em
}

\setlist[enumerate]{
  leftmargin=2em,
  itemsep=0.35em,
  topsep=0.35em
}


% ================================================================
% NUMBERING
% ================================================================

\numberwithin{equation}{chapter}


% ================================================================
% COMMON NUMBER SYSTEMS
% ================================================================

\newcommand{\R}{\mathbb R}
\newcommand{\Rone}{\mathbb R^1}
\newcommand{\Rtwo}{\mathbb R^2}
\newcommand{\Rthree}{\mathbb R^3}
\newcommand{\Rn}{\mathbb R^n}
\newcommand{\Rm}{\mathbb R^m}

\newcommand{\N}{\mathbb N}
\newcommand{\Z}{\mathbb Z}
\newcommand{\Q}{\mathbb Q}
\newcommand{\C}{\mathbb C}


% ================================================================
% SET NOTATION
% ================================================================

% Example:
%   \Set{x \in \R}{x > 0}
%
% Produces:
%   { x in R : x > 0 }

\newcommand{\Set}[2]{\left\{\, #1 \;:\; #2 \,\right\}}

% Alternative with vertical bar:
%   \Setbar{x \in \R}{x > 0}

\newcommand{\Setbar}[2]{\left\{\, #1 \;\middle|\; #2 \,\right\}}

\newcommand{\st}{\;:\;}
\newcommand{\suchthat}{\;\middle|\;}


% ================================================================
% VECTOR NOTATION
% ================================================================
%
% We use arrow notation throughout:
%
%   \vec v, \vec u, \vec x
%
% We do not use boldface for vectors.
% The standard LaTeX \vec command is left unchanged.

\newcommand{\vzero}{\vec 0}

\newcommand{\va}{\vec a}
\newcommand{\vb}{\vec b}
\newcommand{\vc}{\vec c}
\newcommand{\vu}{\vec u}
\newcommand{\vv}{\vec v}
\newcommand{\vw}{\vec w}
\newcommand{\vx}{\vec x}
\newcommand{\vy}{\vec y}
\newcommand{\vz}{\vec z}
\newcommand{\vp}{\vec p}
\newcommand{\vq}{\vec q}

% Standard basis vectors.

\newcommand{\eone}{\vec e_1}
\newcommand{\etwo}{\vec e_2}
\newcommand{\ethree}{\vec e_3}
\newcommand{\ei}{\vec e_i}
\newcommand{\ej}{\vec e_j}

% Optional physics/engineering notation.
% Main notes should primarily use \vec e_1, \vec e_2, \vec e_3.

\newcommand{\ihat}{\hat{\imath}}
\newcommand{\jhat}{\hat{\jmath}}
\newcommand{\khat}{\hat{k}}


% ================================================================
% COLUMN VECTORS AND MATRICES
% ================================================================
%
% We emphasize column vectors.
%
% Examples:
%
%   \twovec{2}{3}
%   \threevec{1}{0}{-4}
%   \colvec{x_1 \\ x_2 \\ \vdots \\ x_n}
%   \mat{1 & 2 \\ 3 & 4}

\newcommand{\colvec}[1]{\begin{bmatrix} #1 \end{bmatrix}}

\newcommand{\twovec}[2]{
\begin{bmatrix}
#1\\
#2
\end{bmatrix}
}

\newcommand{\threevec}[3]{
\begin{bmatrix}
#1\\
#2\\
#3
\end{bmatrix}
}

\newcommand{\fourvec}[4]{
\begin{bmatrix}
#1\\
#2\\
#3\\
#4
\end{bmatrix}
}

\newcommand{\mat}[1]{\begin{bmatrix} #1 \end{bmatrix}}

\newcommand{\smallmat}[1]{%
\left[\begin{smallmatrix} #1 \end{smallmatrix}\right]%
}

% Augmented matrix.
%
% Example:
%
% \augmat{cc|c}{
% 1 & 2 & 3\\
% 4 & 5 & 6
% }

\newcommand{\augmat}[2]{
\left[
\begin{array}{#1}
#2
\end{array}
\right]
}


% ================================================================
% INNER PRODUCTS, NORMS, DISTANCE
% ================================================================

\newcommand{\dotp}[2]{#1 \cdot #2}
\newcommand{\inner}[2]{\left\langle #1, #2 \right\rangle}

\newcommand{\norm}[1]{\left\lVert #1 \right\rVert}
\newcommand{\abs}[1]{\left| #1 \right|}
\newcommand{\dist}[2]{\operatorname{dist}\!\left(#1,#2\right)}


% ================================================================
% LINEAR ALGEBRA OPERATORS
% ================================================================

\DeclareMathOperator{\Span}{span}
\DeclareMathOperator{\Col}{Col}
\DeclareMathOperator{\Nul}{Nul}
\DeclareMathOperator{\Row}{Row}
\DeclareMathOperator{\rank}{rank}
\DeclareMathOperator{\nullity}{nullity}
\DeclareMathOperator{\im}{im}
\DeclareMathOperator{\range}{range}
\DeclareMathOperator{\projop}{proj}
\DeclareMathOperator{\diag}{diag}
\DeclareMathOperator{\tr}{tr}

% Projection notation.
%
% Example:
%   \proj{\vec u}{\vec v}
%
% Means:
%   projection of \vec v onto the direction/subspace \vec u.

\newcommand{\proj}[2]{\projop_{#1}\!\left(#2\right)}

% Matrix transpose and inverse.
%
% Examples:
%   A\T
%   A\inv

\newcommand{\T}{^{\mathsf T}}
\newcommand{\inv}{^{-1}}

% Coordinates relative to a basis.
%
% Example:
%   \coords{\vec v}{B}
%
% Produces:
%   [v]_B

\newcommand{\coords}[2]{\left[#1\right]_{#2}}


% ================================================================
% TEXT EMPHASIS MACROS
% ================================================================

\newcommand{\defterm}[1]{\textbf{#1}}

\newcommand{\important}[1]{%
\begin{center}
\fbox{\parbox{0.85\textwidth}{#1}}
\end{center}
}

\newcommand{\bigidea}[1]{%
\begin{center}
\begin{tcolorbox}[
  width=0.88\textwidth,
  colback=yellow!8,
  colframe=yellow!50!black,
  boxrule=0.8pt,
  arc=2mm
]
#1
\end{tcolorbox}
\end{center}
}


% ================================================================
% THEOREM-LIKE ENVIRONMENTS
% ================================================================

\theoremstyle{definition}

\newtheorem{definition}{Definition}[chapter]
\newtheorem{example}{Example}[chapter]
\newtheorem{exerciseinner}{Exercise}[chapter]

\theoremstyle{plain}

\newtheorem{theorem}{Theorem}[chapter]
\newtheorem{proposition}{Proposition}[chapter]
\newtheorem{fact}{Fact}[chapter]

\theoremstyle{remark}

\newtheorem{remark}{Remark}[chapter]


% ================================================================
% BOX STYLING FOR THEOREM-LIKE ENVIRONMENTS
% ================================================================

\tcolorboxenvironment{definition}{
  enhanced,
  breakable,
  colback=blue!3,
  colframe=blue!55!black,
  boxrule=0.7pt,
  arc=2mm,
  left=1mm,
  right=1mm,
  top=1mm,
  bottom=1mm
}

\tcolorboxenvironment{example}{
  enhanced,
  breakable,
  colback=green!3,
  colframe=green!45!black,
  boxrule=0.7pt,
  arc=2mm,
  left=1mm,
  right=1mm,
  top=1mm,
  bottom=1mm
}

\tcolorboxenvironment{theorem}{
  enhanced,
  breakable,
  colback=purple!3,
  colframe=purple!55!black,
  boxrule=0.7pt,
  arc=2mm,
  left=1mm,
  right=1mm,
  top=1mm,
  bottom=1mm
}

\tcolorboxenvironment{proposition}{
  enhanced,
  breakable,
  colback=purple!3,
  colframe=purple!45!black,
  boxrule=0.7pt,
  arc=2mm,
  left=1mm,
  right=1mm,
  top=1mm,
  bottom=1mm
}

\tcolorboxenvironment{fact}{
  enhanced,
  breakable,
  colback=orange!5,
  colframe=orange!60!black,
  boxrule=0.7pt,
  arc=2mm,
  left=1mm,
  right=1mm,
  top=1mm,
  bottom=1mm
}

\tcolorboxenvironment{remark}{
  enhanced,
  breakable,
  colback=gray!5,
  colframe=gray!50!black,
  boxrule=0.5pt,
  arc=2mm,
  left=1mm,
  right=1mm,
  top=1mm,
  bottom=1mm
}


% ================================================================
% EXERCISES AND SOLUTIONS
% ================================================================
%
% Exercise:
%
% \begin{exercise}
% ...
% \end{exercise}
%
% Exercise with title:
%
% \begin{exercise}[Interpreting a line]
% ...
% \end{exercise}

\NewDocumentEnvironment{exercise}{O{}}
{
  \ifstrempty{#1}
    {\begin{exerciseinner}}
    {\begin{exerciseinner}[#1]}
}
{
  \end{exerciseinner}
}

% Solutions are hidden by default.
% To show solutions, change \showsolutionstrue below.

\newif\ifshowsolutions
\showsolutionsfalse
% \showsolutionstrue

\ifshowsolutions
  \NewDocumentEnvironment{solution}{}
  {
    \begin{proof}[Solution]
  }
  {
    \end{proof}
  }
\else
  \excludecomment{solution}
\fi


% ================================================================
% CHAPTER-LEVEL STRUCTURE BOXES
% ================================================================

\newenvironment{chapteroverview}
{
  \begin{tcolorbox}[
    enhanced,
    breakable,
    colback=cyan!4,
    colframe=cyan!50!black,
    title=\textbf{Chapter Overview},
    boxrule=0.7pt,
    arc=2mm
  ]
}
{
  \end{tcolorbox}
}

\newenvironment{learningobjectives}
{
  \begin{tcolorbox}[
    enhanced,
    breakable,
    colback=cyan!4,
    colframe=cyan!50!black,
    title=\textbf{Learning Objectives},
    boxrule=0.7pt,
    arc=2mm
  ]
  \begin{itemize}
}
{
  \end{itemize}
  \end{tcolorbox}
}

\newenvironment{summarybox}
{
  \begin{tcolorbox}[
    enhanced,
    breakable,
    colback=yellow!5,
    colframe=yellow!50!black,
    title=\textbf{Summary},
    boxrule=0.7pt,
    arc=2mm
  ]
  \begin{itemize}
}
{
  \end{itemize}
  \end{tcolorbox}
}

\newenvironment{checkyourunderstanding}
{
  \begin{tcolorbox}[
    enhanced,
    breakable,
    colback=red!3,
    colframe=red!50!black,
    title=\textbf{Check Your Understanding},
    boxrule=0.7pt,
    arc=2mm
  ]
  \begin{enumerate}
}
{
  \end{enumerate}
  \end{tcolorbox}
}


% ================================================================
% OPTIONAL SECTION-LEVEL BOXES
% ================================================================

\newenvironment{motivation}
{
  \begin{tcolorbox}[
    enhanced,
    breakable,
    colback=green!4,
    colframe=green!45!black,
    title=\textbf{Motivation},
    boxrule=0.7pt,
    arc=2mm
  ]
}
{
  \end{tcolorbox}
}

\newenvironment{warningbox}
{
  \begin{tcolorbox}[
    enhanced,
    breakable,
    colback=red!4,
    colframe=red!60!black,
    title=\textbf{Warning},
    boxrule=0.7pt,
    arc=2mm
  ]
}
{
  \end{tcolorbox}
}

\newenvironment{computationalnote}
{
  \begin{tcolorbox}[
    enhanced,
    breakable,
    colback=gray!6,
    colframe=gray!60!black,
    title=\textbf{Computational Note},
    boxrule=0.7pt,
    arc=2mm
  ]
}
{
  \end{tcolorbox}
}


% ================================================================
% TIKZ STYLES
% ================================================================

\tikzset{
  vector/.style={
    -{Latex[length=3mm]},
    thick,
    blue!70!black
  },
  axis/.style={
    -{Latex[length=2.5mm]},
    gray!70
  },
  point/.style={
    circle,
    fill=black,
    inner sep=1.5pt
  },
  dashedline/.style={
    dashed,
    gray!70
  }
}

% Simple coordinate axes for \R^2 sketches.
%
% Usage inside tikzpicture:
%
% \drawRtwoaxes{-3}{3}{-2}{2}
%
% Arguments:
%   xmin, xmax, ymin, ymax

\NewDocumentCommand{\drawRtwoaxes}{m m m m}
{
  \draw[axis] (#1,0) -- (#2,0) node[right] {$x$};
  \draw[axis] (0,#3) -- (0,#4) node[above] {$y$};
}


% ================================================================
% CLEVEREF NAMES
% ================================================================

\crefname{definition}{Definition}{Definitions}
\Crefname{definition}{Definition}{Definitions}

\crefname{example}{Example}{Examples}
\Crefname{example}{Example}{Examples}

\crefname{theorem}{Theorem}{Theorems}
\Crefname{theorem}{Theorem}{Theorems}

\crefname{proposition}{Proposition}{Propositions}
\Crefname{proposition}{Proposition}{Propositions}

\crefname{fact}{Fact}{Facts}
\Crefname{fact}{Fact}{Facts}

\crefname{remark}{Remark}{Remarks}
\Crefname{remark}{Remark}{Remarks}

\crefname{exerciseinner}{Exercise}{Exercises}
\Crefname{exerciseinner}{Exercise}{Exercises}


% ================================================================
% END OF PREAMBLE
% ================================================================

%
% \begin{document}
% \frontmatter
% \maketitle
% \tableofcontents
%
% \mainmatter
% \chapter{Introduction}
% ...
%
% \end{document}
%
% ================================================================


% ================================================================
% BASIC PACKAGES
% ================================================================

\usepackage[utf8]{inputenc}
\usepackage[T1]{fontenc}
\usepackage{lmodern}
\usepackage{microtype}

\usepackage[margin=1in]{geometry}

\usepackage{amsmath, amssymb, amsthm, mathtools}
\usepackage{graphicx}
\usepackage{xcolor}
\usepackage{booktabs}
\usepackage{array}
\usepackage{enumitem}
\usepackage{comment}
\usepackage{xparse}
\usepackage{etoolbox}


% ================================================================
% FIGURES AND DRAWINGS
% ================================================================

\usepackage{tikz}

\usetikzlibrary{
  arrows.meta,
  calc,
  intersections,
  positioning,
  decorations.pathreplacing,
  angles,
  quotes
}



\usepackage{pgfplots}
\pgfplotsset{compat=1.18}

\definecolor{notespink}{RGB}{210,45,105}
\definecolor{notespurple}{RGB}{95,45,170}


% ================================================================
% BOXES
% ================================================================

\usepackage[most]{tcolorbox}


% ================================================================
% HYPERLINKS AND REFERENCES
% ================================================================

\usepackage{hyperref}

\hypersetup{
  colorlinks=true,
  linkcolor=blue!60!black,
  citecolor=blue!60!black,
  urlcolor=blue!70!black,
  pdftitle={Linear Algebra},
  pdfauthor={},
  pdfsubject={Linear Algebra},
  pdfkeywords={linear algebra, vectors, matrices}
}

\usepackage[nameinlink]{cleveref}


% ================================================================
% PAGE STYLE
% ================================================================

\usepackage{fancyhdr}

\pagestyle{fancy}
\fancyhf{}
\fancyhead[L]{Linear Algebra}
\fancyhead[R]{\leftmark}
\fancyfoot[C]{\thepage}

\setlength{\headheight}{14pt}


% ================================================================
% PARAGRAPH AND LIST FORMATTING
% ================================================================

\setlength{\parindent}{0pt}
\setlength{\parskip}{0.65em}

\setlist[itemize]{
  leftmargin=2em,
  itemsep=0.25em,
  topsep=0.25em
}

\setlist[enumerate]{
  leftmargin=2em,
  itemsep=0.35em,
  topsep=0.35em
}


% ================================================================
% NUMBERING
% ================================================================

\numberwithin{equation}{chapter}


% ================================================================
% COMMON NUMBER SYSTEMS
% ================================================================

\newcommand{\R}{\mathbb R}
\newcommand{\Rone}{\mathbb R^1}
\newcommand{\Rtwo}{\mathbb R^2}
\newcommand{\Rthree}{\mathbb R^3}
\newcommand{\Rn}{\mathbb R^n}
\newcommand{\Rm}{\mathbb R^m}

\newcommand{\N}{\mathbb N}
\newcommand{\Z}{\mathbb Z}
\newcommand{\Q}{\mathbb Q}
\newcommand{\C}{\mathbb C}


% ================================================================
% SET NOTATION
% ================================================================

% Example:
%   \Set{x \in \R}{x > 0}
%
% Produces:
%   { x in R : x > 0 }

\newcommand{\Set}[2]{\left\{\, #1 \;:\; #2 \,\right\}}

% Alternative with vertical bar:
%   \Setbar{x \in \R}{x > 0}

\newcommand{\Setbar}[2]{\left\{\, #1 \;\middle|\; #2 \,\right\}}

\newcommand{\st}{\;:\;}
\newcommand{\suchthat}{\;\middle|\;}


% ================================================================
% VECTOR NOTATION
% ================================================================
%
% We use arrow notation throughout:
%
%   \vec v, \vec u, \vec x
%
% We do not use boldface for vectors.
% The standard LaTeX \vec command is left unchanged.

\newcommand{\vzero}{\vec 0}

\newcommand{\va}{\vec a}
\newcommand{\vb}{\vec b}
\newcommand{\vc}{\vec c}
\newcommand{\vu}{\vec u}
\newcommand{\vv}{\vec v}
\newcommand{\vw}{\vec w}
\newcommand{\vx}{\vec x}
\newcommand{\vy}{\vec y}
\newcommand{\vz}{\vec z}
\newcommand{\vp}{\vec p}
\newcommand{\vq}{\vec q}

% Standard basis vectors.

\newcommand{\eone}{\vec e_1}
\newcommand{\etwo}{\vec e_2}
\newcommand{\ethree}{\vec e_3}
\newcommand{\ei}{\vec e_i}
\newcommand{\ej}{\vec e_j}

% Optional physics/engineering notation.
% Main notes should primarily use \vec e_1, \vec e_2, \vec e_3.

\newcommand{\ihat}{\hat{\imath}}
\newcommand{\jhat}{\hat{\jmath}}
\newcommand{\khat}{\hat{k}}


% ================================================================
% COLUMN VECTORS AND MATRICES
% ================================================================
%
% We emphasize column vectors.
%
% Examples:
%
%   \twovec{2}{3}
%   \threevec{1}{0}{-4}
%   \colvec{x_1 \\ x_2 \\ \vdots \\ x_n}
%   \mat{1 & 2 \\ 3 & 4}

\newcommand{\colvec}[1]{\begin{bmatrix} #1 \end{bmatrix}}

\newcommand{\twovec}[2]{
\begin{bmatrix}
#1\\
#2
\end{bmatrix}
}

\newcommand{\threevec}[3]{
\begin{bmatrix}
#1\\
#2\\
#3
\end{bmatrix}
}

\newcommand{\fourvec}[4]{
\begin{bmatrix}
#1\\
#2\\
#3\\
#4
\end{bmatrix}
}

\newcommand{\mat}[1]{\begin{bmatrix} #1 \end{bmatrix}}

\newcommand{\smallmat}[1]{%
\left[\begin{smallmatrix} #1 \end{smallmatrix}\right]%
}

% Augmented matrix.
%
% Example:
%
% \augmat{cc|c}{
% 1 & 2 & 3\\
% 4 & 5 & 6
% }

\newcommand{\augmat}[2]{
\left[
\begin{array}{#1}
#2
\end{array}
\right]
}


% ================================================================
% INNER PRODUCTS, NORMS, DISTANCE
% ================================================================

\newcommand{\dotp}[2]{#1 \cdot #2}
\newcommand{\inner}[2]{\left\langle #1, #2 \right\rangle}

\newcommand{\norm}[1]{\left\lVert #1 \right\rVert}
\newcommand{\abs}[1]{\left| #1 \right|}
\newcommand{\dist}[2]{\operatorname{dist}\!\left(#1,#2\right)}


% ================================================================
% LINEAR ALGEBRA OPERATORS
% ================================================================

\DeclareMathOperator{\Span}{span}
\DeclareMathOperator{\Col}{Col}
\DeclareMathOperator{\Nul}{Nul}
\DeclareMathOperator{\Row}{Row}
\DeclareMathOperator{\rank}{rank}
\DeclareMathOperator{\nullity}{nullity}
\DeclareMathOperator{\im}{im}
\DeclareMathOperator{\range}{range}
\DeclareMathOperator{\projop}{proj}
\DeclareMathOperator{\diag}{diag}
\DeclareMathOperator{\tr}{tr}

% Projection notation.
%
% Example:
%   \proj{\vec u}{\vec v}
%
% Means:
%   projection of \vec v onto the direction/subspace \vec u.

\newcommand{\proj}[2]{\projop_{#1}\!\left(#2\right)}

% Matrix transpose and inverse.
%
% Examples:
%   A\T
%   A\inv

\newcommand{\T}{^{\mathsf T}}
\newcommand{\inv}{^{-1}}

% Coordinates relative to a basis.
%
% Example:
%   \coords{\vec v}{B}
%
% Produces:
%   [v]_B

\newcommand{\coords}[2]{\left[#1\right]_{#2}}


% ================================================================
% TEXT EMPHASIS MACROS
% ================================================================

\newcommand{\defterm}[1]{\textbf{#1}}

\newcommand{\important}[1]{%
\begin{center}
\fbox{\parbox{0.85\textwidth}{#1}}
\end{center}
}

\newcommand{\bigidea}[1]{%
\begin{center}
\begin{tcolorbox}[
  width=0.88\textwidth,
  colback=yellow!8,
  colframe=yellow!50!black,
  boxrule=0.8pt,
  arc=2mm
]
#1
\end{tcolorbox}
\end{center}
}


% ================================================================
% THEOREM-LIKE ENVIRONMENTS
% ================================================================

\theoremstyle{definition}

\newtheorem{definition}{Definition}[chapter]
\newtheorem{example}{Example}[chapter]
\newtheorem{exerciseinner}{Exercise}[chapter]

\theoremstyle{plain}

\newtheorem{theorem}{Theorem}[chapter]
\newtheorem{proposition}{Proposition}[chapter]
\newtheorem{fact}{Fact}[chapter]

\theoremstyle{remark}

\newtheorem{remark}{Remark}[chapter]


% ================================================================
% BOX STYLING FOR THEOREM-LIKE ENVIRONMENTS
% ================================================================

\tcolorboxenvironment{definition}{
  enhanced,
  breakable,
  colback=blue!3,
  colframe=blue!55!black,
  boxrule=0.7pt,
  arc=2mm,
  left=1mm,
  right=1mm,
  top=1mm,
  bottom=1mm
}

\tcolorboxenvironment{example}{
  enhanced,
  breakable,
  colback=green!3,
  colframe=green!45!black,
  boxrule=0.7pt,
  arc=2mm,
  left=1mm,
  right=1mm,
  top=1mm,
  bottom=1mm
}

\tcolorboxenvironment{theorem}{
  enhanced,
  breakable,
  colback=purple!3,
  colframe=purple!55!black,
  boxrule=0.7pt,
  arc=2mm,
  left=1mm,
  right=1mm,
  top=1mm,
  bottom=1mm
}

\tcolorboxenvironment{proposition}{
  enhanced,
  breakable,
  colback=purple!3,
  colframe=purple!45!black,
  boxrule=0.7pt,
  arc=2mm,
  left=1mm,
  right=1mm,
  top=1mm,
  bottom=1mm
}

\tcolorboxenvironment{fact}{
  enhanced,
  breakable,
  colback=orange!5,
  colframe=orange!60!black,
  boxrule=0.7pt,
  arc=2mm,
  left=1mm,
  right=1mm,
  top=1mm,
  bottom=1mm
}

\tcolorboxenvironment{remark}{
  enhanced,
  breakable,
  colback=gray!5,
  colframe=gray!50!black,
  boxrule=0.5pt,
  arc=2mm,
  left=1mm,
  right=1mm,
  top=1mm,
  bottom=1mm
}


% ================================================================
% EXERCISES AND SOLUTIONS
% ================================================================
%
% Exercise:
%
% \begin{exercise}
% ...
% \end{exercise}
%
% Exercise with title:
%
% \begin{exercise}[Interpreting a line]
% ...
% \end{exercise}

\NewDocumentEnvironment{exercise}{O{}}
{
  \ifstrempty{#1}
    {\begin{exerciseinner}}
    {\begin{exerciseinner}[#1]}
}
{
  \end{exerciseinner}
}

% Solutions are hidden by default.
% To show solutions, change \showsolutionstrue below.

\newif\ifshowsolutions
\showsolutionsfalse
% \showsolutionstrue

\ifshowsolutions
  \NewDocumentEnvironment{solution}{}
  {
    \begin{proof}[Solution]
  }
  {
    \end{proof}
  }
\else
  \excludecomment{solution}
\fi


% ================================================================
% CHAPTER-LEVEL STRUCTURE BOXES
% ================================================================

\newenvironment{chapteroverview}
{
  \begin{tcolorbox}[
    enhanced,
    breakable,
    colback=cyan!4,
    colframe=cyan!50!black,
    title=\textbf{Chapter Overview},
    boxrule=0.7pt,
    arc=2mm
  ]
}
{
  \end{tcolorbox}
}

\newenvironment{learningobjectives}
{
  \begin{tcolorbox}[
    enhanced,
    breakable,
    colback=cyan!4,
    colframe=cyan!50!black,
    title=\textbf{Learning Objectives},
    boxrule=0.7pt,
    arc=2mm
  ]
  \begin{itemize}
}
{
  \end{itemize}
  \end{tcolorbox}
}

\newenvironment{summarybox}
{
  \begin{tcolorbox}[
    enhanced,
    breakable,
    colback=yellow!5,
    colframe=yellow!50!black,
    title=\textbf{Summary},
    boxrule=0.7pt,
    arc=2mm
  ]
  \begin{itemize}
}
{
  \end{itemize}
  \end{tcolorbox}
}

\newenvironment{checkyourunderstanding}
{
  \begin{tcolorbox}[
    enhanced,
    breakable,
    colback=red!3,
    colframe=red!50!black,
    title=\textbf{Check Your Understanding},
    boxrule=0.7pt,
    arc=2mm
  ]
  \begin{enumerate}
}
{
  \end{enumerate}
  \end{tcolorbox}
}


% ================================================================
% OPTIONAL SECTION-LEVEL BOXES
% ================================================================

\newenvironment{motivation}
{
  \begin{tcolorbox}[
    enhanced,
    breakable,
    colback=green!4,
    colframe=green!45!black,
    title=\textbf{Motivation},
    boxrule=0.7pt,
    arc=2mm
  ]
}
{
  \end{tcolorbox}
}

\newenvironment{warningbox}
{
  \begin{tcolorbox}[
    enhanced,
    breakable,
    colback=red!4,
    colframe=red!60!black,
    title=\textbf{Warning},
    boxrule=0.7pt,
    arc=2mm
  ]
}
{
  \end{tcolorbox}
}

\newenvironment{computationalnote}
{
  \begin{tcolorbox}[
    enhanced,
    breakable,
    colback=gray!6,
    colframe=gray!60!black,
    title=\textbf{Computational Note},
    boxrule=0.7pt,
    arc=2mm
  ]
}
{
  \end{tcolorbox}
}


% ================================================================
% TIKZ STYLES
% ================================================================

\tikzset{
  vector/.style={
    -{Latex[length=3mm]},
    thick,
    blue!70!black
  },
  axis/.style={
    -{Latex[length=2.5mm]},
    gray!70
  },
  point/.style={
    circle,
    fill=black,
    inner sep=1.5pt
  },
  dashedline/.style={
    dashed,
    gray!70
  }
}

% Simple coordinate axes for \R^2 sketches.
%
% Usage inside tikzpicture:
%
% \drawRtwoaxes{-3}{3}{-2}{2}
%
% Arguments:
%   xmin, xmax, ymin, ymax

\NewDocumentCommand{\drawRtwoaxes}{m m m m}
{
  \draw[axis] (#1,0) -- (#2,0) node[right] {$x$};
  \draw[axis] (0,#3) -- (0,#4) node[above] {$y$};
}


% ================================================================
% CLEVEREF NAMES
% ================================================================

\crefname{definition}{Definition}{Definitions}
\Crefname{definition}{Definition}{Definitions}

\crefname{example}{Example}{Examples}
\Crefname{example}{Example}{Examples}

\crefname{theorem}{Theorem}{Theorems}
\Crefname{theorem}{Theorem}{Theorems}

\crefname{proposition}{Proposition}{Propositions}
\Crefname{proposition}{Proposition}{Propositions}

\crefname{fact}{Fact}{Facts}
\Crefname{fact}{Fact}{Facts}

\crefname{remark}{Remark}{Remarks}
\Crefname{remark}{Remark}{Remarks}

\crefname{exerciseinner}{Exercise}{Exercises}
\Crefname{exerciseinner}{Exercise}{Exercises}


% ================================================================
% END OF PREAMBLE
% ================================================================
